% Part: applied-modal-logics
% Chapter: temporal-logic
% Section: properties-accessibility

\documentclass[../../../include/open-logic-section]{subfiles}

\begin{document}

\olfileid{aml}{tl}{acc}

\olsection{时态框架的性质}

由于我们的时态模型没有对关系~$\prec$ 施加任何条件，
在我们熟悉的公理中，在所有模型中都成立的仅有~$K$，
或其类似物 $K_{\Gtemp}$ 和~$K_{\Htemp}$：\begin{align*}
\tag{$K_{\Gtemp}$}  \Gtemp (p \to q) & \to (\Gtemp p \to \Gtemp q)\\
\tag{$K_{\Htemp}$}  \Htemp (p \to q) & \to (\Htemp p \to \Htemp q)
\end{align*}

然而，如果我们想对模型表示时间的方式施加更严格的条件，
例如确保 $\prec$ 是一个线序，
那么在我们的模型中就会产生其他的有效性。

\begin{table}[t]
  \begin{tabular}{| p{.48\textwidth} || p{.48\textwidth} |}
    \hline
    {\emph{若 $\prec$ 是……}} & {\emph{则……在~$\mModel{M}$ 中为真：}} \\
    \hline \hline
    \emph{传递的}： \newline
    $\forall u \forall v \forall w ((u \prec v \land v \prec w) \lif u \prec w)$ &
    $\Ftemp \Ftemp p \lif \Ftemp p$  \\
    \hline
    \emph{线性的}： \newline
    $\forall w \forall v (w \prec v \lor w = v \lor v \prec w)$ &
    $(\Ftemp \Ptemp  p \lor \Ptemp \Ftemp p) \lif (\Ptemp p \lor p \lor \Ftemp p)$\\
    \hline
    \emph{稠密的}： \newline
    $\forall w \forall v (w \prec v \to \exists u(w \prec u \land u \prec v))$ &
    $\Ftemp p \lif \Ftemp \Ftemp p$ \\
    \hline
    \emph{无界的（过去）}： \newline
    $\forall w \exists v( v \prec w)$ &
    $\Htemp p \to \Ptemp p$ \\
    \hline
    \emph{无界的（未来）}： \newline
    $\forall w \exists v( w \prec v)$ &
    $\Gtemp p \to \Ftemp p$ \\
    \hline
  \end{tabular}
  \caption{时态框架的一些对应性质。}
  \ollabel{tab:correspondence}
\end{table}

对于旨在表示时间的模型而言，
\olref{tab:correspondence} 中的一些性质似乎是理想的特征。
然而，值得注意的是，尽管我们可以对 $\prec$ 关系施加任意条件，
但并非所有条件都对应着能在时态逻辑语言中表达的公式。
例如，非自反性（即 $\forall w \lnot (w \prec w)$），
在时态逻辑中就没有对应的公式。

\end{document}